\documentclass{article}

\usepackage{amsmath,amssymb}
\usepackage{xurl}
\usepackage[hidelinks]{hyperref}

% Shared commands for notes in docs/paper-gaps/.

\DeclareUnicodeCharacter{2080}{\ifmmode{}_{0}\else\textsubscript{0}\fi}
\DeclareUnicodeCharacter{2081}{\ifmmode{}_{1}\else\textsubscript{1}\fi}
\DeclareUnicodeCharacter{2082}{\ifmmode{}_{2}\else\textsubscript{2}\fi}
\DeclareUnicodeCharacter{2099}{\ifmmode{}_{n}\else\textsubscript{n}\fi}

\newcommand{\Lean}{\textsc{Lean}}
\ExplSyntaxOn
\NewDocumentCommand{\leanid}{m}{
  \ifmmode
    \textsf{\detokenize{#1}}
  \else
    \group_begin:
    \urlstyle{sf}
    \tl_set:Nn \l_tmpa_tl {#1}
    \tl_replace_all:Nnn \l_tmpa_tl {\_} {_}
    \exp_args:NV \path \l_tmpa_tl
    \group_end:
  \fi
}
\ExplSyntaxOff
\providecommand{\tr}{\operatorname{tr}}
\newcommand{\tnleanIssueBase}{https://github.com/LionSR/TNLean/issues/}
\newcommand{\tnleanPrBase}{https://github.com/LionSR/TNLean/pull/}
\newcommand{\tnleanDocBase}{https://sirui-lu.com/TNLean/docs/}
\newcommand{\ghissue}[1]{\href{\tnleanIssueBase#1}{GitHub issue \##1}}
\newcommand{\ghpr}[1]{\href{\tnleanPrBase#1}{PR \##1}}
\newcommand{\doclink}[1]{\href{\tnleanDocBase#1.html}{\path{#1}}}

% Dirac notation and the identity operator, as in blueprint/src/macros/common.tex.
% \providecommand keeps note-local definitions authoritative.
\usepackage{dsfont}
\providecommand{\ket}[1]{|#1\rangle}
\providecommand{\bra}[1]{\langle #1|}
\providecommand{\braket}[2]{\langle #1 | #2 \rangle}
\providecommand{\ketbra}[2]{|#1\rangle\!\langle #2|}
\providecommand{\Id}{\mathds{1}}
\providecommand{\one}{\mathds{1}}
\providecommand{\C}{\mathbb{C}}
\providecommand{\MN}[1]{M_{#1}(\C)}
\providecommand{\MD}{M_D(\C)}

% External referencing for paper sources.  The build helper writes sanitized
% auxiliary files under build/paper-aux/ and excludes duplicated source labels.
\usepackage{xr}

% Import a generated paper auxiliary file with a caller-supplied label prefix.
% The candidate paths support builds from docs/paper-gaps/, the repository root,
% and build/paper-aux/ itself.
\newcommand{\importpaperaux}[2]{%
  \IfFileExists{../../build/paper-aux/#2.aux}{%
    \externaldocument[#1]{../../build/paper-aux/#2}%
  }{%
    \IfFileExists{build/paper-aux/#2.aux}{%
      \externaldocument[#1]{build/paper-aux/#2}%
    }{%
      \IfFileExists{#2.aux}{%
        \externaldocument[#1]{#2}%
      }{}%
    }%
  }%
}

% General external reference macros
\newcommand{\paperref}[2]{\ref{#1-#2}}
\newcommand{\papereqref}[2]{(\ref{#1-#2})}

% CPSV16 source (1606.00608 / MPDO-22-12-17-2.tex) external document and shortcuts
\importpaperaux{cpsv16-}{1606.00608/MPDO-22-12-17-2-tnlean}
\newcommand{\cpsvref}[1]{\paperref{cpsv16}{#1}}
\newcommand{\cpsveqref}[1]{\papereqref{cpsv16}{#1}}

% Verdict marker: \gapnote{<kind>}{<status>}. Kind is the policy
% classification (clarification, local-correction, scope-restriction,
% unfaithful, false-source, open-gap); status is open, wip (elimination
% underway), resolved, or historical. Machine-read by the site index;
% prints a verdict line.
\newcommand{\gapnote}[2]{\par\noindent\textbf{Verdict:} #1 (#2).\par}


\title{The Index Typo in the Pure-State Renormalization-Flow Equation}
\author{}
\date{2026-07-30}

\begin{document}
\maketitle
\gapnote{local-correction}{resolved}

\section{Malformed source display}

Theorem~3.1 of Ref.~\cite{Cirac2016MPDO_arXiv} characterizes limits of the
pure-state renormalization flow. Its equation at source lines~400--403 is
printed as
\begin{align}
    A^{i_1}A^{i_2}
        &=
        \sum_{i_1,i_2,j}
        U_{(i_1,i_2),j}A^{j_1}.
    \label{eq:cpsv16_renormalization_flow_index_typo_malformed_source}
\end{align}
The indices $i_1,i_2$ are free on the left-hand side of
Eq.~\ref{eq:cpsv16_renormalization_flow_index_typo_malformed_source} but are
summed on its right-hand side. The index $j_1$ is not introduced. The printed
display is therefore not a well-formed component identity.

\section{Intended equation}

The renormalization step at source lines~389--394 of
Ref.~\cite{Cirac2016MPDO_arXiv} first blocks two physical indices and then
writes the blocked matrices as a physical isometry applied to a one-site
family. At a fixed point with a common physical dimension, the component
identity is
\begin{align}
    A^{i_1}A^{i_2}
        &=
        \sum_j U_{(i_1,i_2),j}A^j,
        \qquad\text{for every }i_1,i_2.
    \label{eq:cpsv16_renormalization_flow_index_typo_corrected_blocking}
\end{align}
The diagrams at source lines~407--410 of
Ref.~\cite{Cirac2016MPDO_arXiv} express
Eq.~\ref{eq:cpsv16_renormalization_flow_index_typo_corrected_blocking}: the two
external physical legs $i_1,i_2$ enter $U$, and only the one-site index $j$ is
summed.

The Appendix proof at source lines~1205--1209 of
Ref.~\cite{Cirac2016MPDO_arXiv} reaches the same identity. It first identifies
the completely positive maps by
\begin{align}
    \mathcal E^2 &= \mathcal E,
    \label{eq:cpsv16_renormalization_flow_index_typo_cp_idempotence}
\end{align}
and then relates the corresponding Kraus families by a physical isometry,
which gives
Eq.~\ref{eq:cpsv16_renormalization_flow_index_typo_corrected_blocking}.

\section{Formal treatment}

The formal theorem uses
Eq.~\ref{eq:cpsv16_renormalization_flow_index_typo_corrected_blocking}. This is
a local correction to the summation and output indices in the printed display
of Ref.~\cite{Cirac2016MPDO_arXiv}. It changes neither the hypothesis that the
tensor is a limit of the renormalization flow nor the conclusion that a
physical isometry relates the blocked tensor to the original tensor.

\bibliographystyle{alpha}
\bibliography{references}

\end{document}
